\documentclass[10pt]{article}

\usepackage[margin=0.6in]{geometry}
\usepackage[T1]{fontenc}
\usepackage[utf8]{inputenc}
\usepackage{lmodern}
\usepackage{microtype}
\usepackage{booktabs}
\usepackage{longtable}
\usepackage{tabularx}
\usepackage{array}
\usepackage{pdflscape}
\usepackage{xcolor}
\usepackage{hyperref}
\usepackage{float}

\hypersetup{colorlinks=true,linkcolor=blue!60!black,urlcolor=blue!60!black}

\newcolumntype{Y}{>{\raggedright\arraybackslash}X}
\newcommand{\code}[1]{\texttt{#1}}

\title{Results-Only Comparison of Models and Prompt Variants\\Using \code{resource/results\_variants.txt}}
\author{}
\date{April 1, 2026}

\begin{document}
\maketitle

\section*{How to Read This Report}
This report is intentionally results-only. It is based on \code{resource/results\_variants.txt} and focuses on what the logged outputs show under verifier-style criteria:
\begin{itemize}
  \item waypoint-count correctness,
  \item local progress toward the goal,
  \item kinematic plausibility at a qualitative level,
  \item obstacle-aware behavior in the right-obstacle case,
  \item and consistency across prompt variants.
\end{itemize}

Because the logged responses do not include \code{selected\_waypoint\_index}, and some entries are truncated or malformed, the cells below use qualitative judgments instead of strict numeric verifier pass/fail labels.

\section*{Prompt Legend}
\begin{enumerate}
  \item \code{01}: baseline full
  \item \code{02}: no reasoning, keep examples
  \item \code{03}: no reasoning, one example
  \item \code{04}: no reasoning, no examples
  \item \code{05}: reasoning, no examples
  \item \code{06}: concise with reasoning
  \item \code{07}: concise no reasoning
  \item \code{08}: no reasoning, center bias
  \item \code{09}: reasoning, two examples
  \item \code{10}: minimal waypoints only
\end{enumerate}

\section*{Cell Labels}
\begin{itemize}
  \item \textbf{Strong}: five waypoints, locally sensible, clear forward progress, and good task alignment.
  \item \textbf{Good}: mostly sensible, minor issues, but still usable.
  \item \textbf{Mixed}: some progress or correct structure, but notable semantic or stability issues.
  \item \textbf{Weak}: poor local planning quality, wrong directional behavior, malformed output, or clear instability.
  \item \textbf{Fail}: severe structural issue such as missing waypoint count or obviously unusable behavior.
\end{itemize}

\section*{Overall Ranking}
\begin{table}[H]
\centering
\begin{tabularx}{\textwidth}{p{3.2cm}p{2.1cm}p{2.1cm}Y}
\toprule
Model & Row 0 & Row 4 & Summary \\
\midrule
ChatGPT Mini & Strong & Strong & Most consistent model in the log. Stable across all prompt styles with clean 5-waypoint outputs. \\
Qwen 14B & Good--Strong & Good--Strong & Best Qwen model in the file. Usually sensible, with some prompt brittleness and a few regressive variants. \\
Qwen 7B untuned & Mixed & Mixed & Some usable outputs, but high prompt sensitivity and several semantically poor trajectories. \\
Qwen 7B FT post-hoc reasoning & Mixed--Weak & Mixed & Can produce plausible plans, but very inconsistent and often non-local. \\
Qwen 7B FT no reasoning & Weak & Weak & Least reliable family: repeated points, malformed outputs, four-waypoint failures, and unstable motion. \\
\bottomrule
\end{tabularx}
\end{table}

\begin{landscape}
\section*{Row 0: No-Obstacle / Clear-Path Results}
\renewcommand{\arraystretch}{1.25}
\begin{longtable}{p{2.6cm}p{2.35cm}p{2.35cm}p{2.35cm}p{2.35cm}p{2.35cm}p{2.35cm}p{2.35cm}p{2.35cm}p{2.35cm}p{2.35cm}}
\toprule
Model & \code{01} & \code{02} & \code{03} & \code{04} & \code{05} & \code{06} & \code{07} & \code{08} & \code{09} & \code{10} \\
\midrule
\endhead
ChatGPT Mini &
Strong: centered small-step rollout, clean schema, strong verifier-style progress. &
Strong: same stable centerline behavior without reasoning field. &
Good: mild diagonal bias but still local and forward. &
Good: larger steps, but still straight and useful. &
Mixed: geometry is fine, but reasoning text is internally inconsistent with the waypoints. &
Mixed: aggressive long-horizon jumps; still forward, but less local. &
Mixed: same long-horizon compression as 06, less local than ideal. &
Strong: center-bias prompt behaves cleanly and predictably. &
Strong: same strong baseline behavior with reasoning present. &
Strong: minimal prompt still yields neat forward rollout. \\

Qwen 7B untuned &
Good: small diagonal rollout, locally plausible, but reasoning is truncated. &
Good: similar to baseline and structurally clean. &
Weak: starts at the obstacle coordinate; semantically suspicious for a clear-path case. &
Mixed: valid five points, but overly coarse global-style path. &
Weak: waypoints hug current position and reasoning is badly inconsistent/truncated. &
Mixed: diagonal progress is plausible but drifts away from the clear center-progress hint. &
Mixed: forward progress exists, but steps become too coarse and global. &
Good: same usable diagonal baseline pattern. &
Good: same geometry as baseline, but reasoning is truncated. &
Good: short local rollout and decent progress. \\

Qwen 7B FT post-hoc &
Weak: essentially jumps to goal region immediately; not a local rollout. &
Fail: malformed JSON plus near-goal behavior; structurally unreliable. &
Mixed: consistent medium-range diagonal rollout, but not clearly local. &
Weak: starts on obstacle location and then jumps to goal. &
Weak: creeps near obstacle coordinates; poor task alignment. &
Mixed: obstacle-anchored start then diagonal escape; not ideal for clear path. &
Weak: giant non-local jumps in \(y\). &
Weak: immediate jump to goal region with extra overshoot. &
Mixed: near-goal behavior again; better than 01 but still non-local. &
Mixed: monotone near-goal descending points; coherent but not local-planner behavior. \\

Qwen 7B FT no reasoning &
Weak: repeated identical points and truncated reasoning. &
Weak: all waypoints collapse directly to the goal. &
Fail: only four waypoints. &
Strong: simple centered rollout, one of the few clean cases. &
Weak: repeated near-static points and self-contradictory reasoning. &
Good: local straight rollout with mild center bias. &
Good: same stable local rollout without reasoning. &
Weak: again collapses all waypoints to the goal. &
Weak: tiny near-static drift and truncated reasoning. &
Mixed: five points but adds unexpected \(z\) field and repeats current state. \\

Qwen 14B &
Strong: clean, centered, local rollout. &
Strong: same stable centerline behavior. &
Good: diagonal but still locally plausible. &
Good: larger forward spacing, still valid and useful. &
Mixed: waypoints are fine, but reasoning text reports inconsistent progress. &
Mixed: plausible diagonal rollout, but drifts more than needed in the clear case. &
Mixed: very long forward jumps, less local than baseline. &
Strong: same clean baseline response. &
Good: similar to untuned 7B diagonal rollout; usable but truncated reasoning. &
Good: minimal prompt remains coherent, though less centered than baseline. \\
\bottomrule
\end{longtable}

\section*{Row 4: Obstacle on Right / Obstacle-Avoidance Results}
\begin{longtable}{p{2.6cm}p{2.35cm}p{2.35cm}p{2.35cm}p{2.35cm}p{2.35cm}p{2.35cm}p{2.35cm}p{2.35cm}p{2.35cm}p{2.35cm}}
\toprule
Model & \code{01} & \code{02} & \code{03} & \code{04} & \code{05} & \code{06} & \code{07} & \code{08} & \code{09} & \code{10} \\
\midrule
\endhead
ChatGPT Mini &
Strong: forward progress with mild left correction; very plausible. &
Strong: same general behavior, slightly more spread in \(y\). &
Mixed: stays fixed in \(x\) and barely reacts to right-side caution. &
Weak: moves further right, opposite of the left-clearer cue. &
Weak: same rightward drift as 04, plus unreliable reasoning. &
Good: preserves progress and avoids large jumps, though lateral response is modest. &
Good: forward progress is strong, but still somewhat center-locked. &
Mixed: center-bias keeps \(x\) fixed and under-reacts to obstacle cue. &
Good: coherent local rollout, though still conservative in lateral avoidance. &
Strong: minimal prompt gives a clean leftward move toward center. \\

Qwen 7B untuned &
Good: small leftward correction with local progress; reasoning truncated. &
Good: similar to baseline and usable. &
Weak: moves right instead of left, poor obstacle response. &
Weak: keeps moving right and outward, not away from hazard. &
Weak: similar wrong-direction behavior plus inconsistent reasoning. &
Mixed: plausible local spacing, but drifts right instead of left. &
Weak: large forward leaps without convincing obstacle-aware correction. &
Good: small leftward correction again, better than 03--07. &
Good: same as baseline geometry, but reasoning truncated. &
Strong: clear leftward shift toward safer side with local spacing. \\

Qwen 7B FT post-hoc &
Weak: barely moves; almost static around current state. &
Weak: jumps to goal region near \(y \approx 29\), not local. &
Mixed: local diagonal motion, but biased right rather than left. &
Good: small local leftward correction, one of its better row-4 outputs. &
Good: similar to 04, leftward and local. &
Good: same controlled leftward arc with steady spacing. &
Good: similar local leftward correction. &
Weak: mostly local early points, then a huge jump to \(y=29\). &
Strong: best row-4 behavior in this model family; local, leftward, and smooth. &
Good: very local leftward drift, though perhaps too conservative. \\

Qwen 7B FT no reasoning &
Weak: repeated same point, little meaningful motion. &
Mixed: valid five points, but drifts right rather than strongly left. &
Fail: only four waypoints. &
Mixed: local and valid, but too little reaction to the obstacle cue. &
Weak: tiny motion and truncated reasoning; barely responsive. &
Weak: reverses \(y\) direction and moves away from forward progress. &
Weak: same regression in \(y\), with unexpected \(z\) fields. &
Fail: only four waypoints. &
Mixed: steady forward motion, but truncated reasoning and limited avoidance insight. &
Weak: repeated same point, effectively no rollout. \\

Qwen 14B &
Strong: mild but correct leftward avoidance with local progress. &
Strong: similar to baseline and structurally clean. &
Mixed: moves to the right, so obstacle response is not ideal. &
Fail: decreases \(y\), strongly regressive for the task. &
Mixed: forward motion exists, but reasoning claims zero progress and the geometry is too center-locked. &
Strong: strongest leftward avoidance in the log while preserving progress. &
Strong: same as 06, slightly simpler and still very good. &
Strong: balanced mild leftward correction, good compromise between progress and caution. &
Strong: similar to baseline with slightly better leftward drift. &
Strong: minimal prompt still produces a clear leftward avoidance plan. \\
\bottomrule
\end{longtable}
\end{landscape}

\section*{Best Prompt Variants by Model}
\begin{table}[H]
\centering
\begin{tabularx}{\textwidth}{p{3.1cm}p{2.4cm}p{2.4cm}Y}
\toprule
Model & Best row-0 prompts & Best row-4 prompts & Comment \\
\midrule
ChatGPT Mini & \code{01}, \code{02}, \code{08}, \code{09}, \code{10} & \code{01}, \code{02}, \code{10} & Very robust overall; minimal prompt surprisingly strong. \\
Qwen 7B untuned & \code{01}, \code{02}, \code{08}, \code{10} & \code{01}, \code{02}, \code{08}, \code{10} & Needs careful prompt choice; one-example and no-example variants are unreliable. \\
Qwen 7B FT post-hoc & \code{03} only marginally & \code{04}, \code{05}, \code{06}, \code{07}, \code{09}, \code{10} & Better in obstacle case than clear-path case, but still unstable. \\
Qwen 7B FT no reasoning & \code{04}, \code{06}, \code{07} & none clearly strong & Fine-tuning appears to have hurt stability more than helped. \\
Qwen 14B & \code{01}, \code{02}, \code{08} & \code{01}, \code{02}, \code{06}, \code{07}, \code{08}, \code{09}, \code{10} & Best open-model result in the file. \\
\bottomrule
\end{tabularx}
\end{table}

\section*{Takeaways}
\begin{itemize}
  \item ChatGPT Mini is the strongest result in the file because it stays clean and useful across all prompt styles.
  \item Qwen 14B is the best Qwen model and is especially good in the right-obstacle environment when given concise or lightly guided prompts.
  \item Untuned Qwen 7B is sometimes usable, but it is too prompt-sensitive to trust broadly.
  \item The two fine-tuned 7B models are not winning in this log. One is erratic and non-local; the other has repeated-point failures and malformed outputs.
  \item For obstacle avoidance specifically, Qwen 14B with prompts \code{06}, \code{07}, \code{08}, \code{09}, or \code{10} looks strongest among the open models recorded here.
\end{itemize}

\section*{Compile Command}
\begin{verbatim}
cd /home/prachit/Desktop/vlm-conformal/px4_ws/src/llm_drone/resource
pdflatex -interaction=nonstopmode model_prompt_variant_comparison_report.tex
\end{verbatim}

\end{document}
